\section{Módulos clave para construir un AI Agent}

\subsection{Arquitectura de la plataforma Vertex AI}

El diseño de Vertex AI de Google refleja las necesidades comunes de una plataforma de IA empresarial:

\begin{itemize}
    \item \textbf{Capa de chips}: flexibilidad para elegir GPU/TPU (actualmente hay una escasez global severa de chips)
    \item \textbf{Capa de modelos}: una interfaz unificada para modelos propios, de terceros y de código abierto
    \item \textbf{Capa de herramientas}: Model Builder y Agent Builder ofrecen capacidades de personalización y orquestación
    \item \textbf{Capa de aplicaciones}: aplicaciones verticales como búsqueda, conversación y análisis de datos
\end{itemize}

\subsection{Diferenciación de Gemini}

\begin{itemize}
    \item Multimodalidad nativa (mezcla varias modalidades desde el inicio del entrenamiento)
    \item Ventana de contexto extremadamente grande (en los experimentos ya alcanza 10 millones de tokens)
    \item Prueba de Needle-in-a-haystack: recuperación precisa de información en contextos extremadamente largos
\end{itemize}

\subsection{Módulos de Planner y Memory}

Dentro de un Agent normalmente hay tres componentes que colaboran: planner, executor y memory. El flujo típico es:

\begin{itemize}
    \item \textbf{Planner}: descompone la intención del usuario en pasos (por ejemplo, búsqueda, llamada a herramientas, resumen)
    \item \textbf{Memory}: registra cada conversación, consulta y salida de herramientas, y ofrece context re-use
    \item \textbf{Executor}: se encarga de invocar herramientas específicas, API o entornos de ejecución personalizados
\end{itemize}

\begin{importantbox}{El reto de sincronización entre planner y memory}
Una vez que el planner define los pasos, memory necesita devolver rápidamente los resultados de ejecuciones anteriores. Si memory no es lo bastante detallada (solo conserva la ronda más reciente), el planner volverá a llamar a las herramientas de forma repetida; si memory es demasiado grande, esto provocará latency. La solución es introducir compressed memory + recall hints, es decir, mantener en memory un query summary como referencia para el planner.
\end{importantbox}
\begin{importantbox}{División de funciones entre Planner y Memory}
Memory le proporciona al planner una caché de contexto; el planner, factores de decisión (el fracaso de la ronda anterior, la tasa de éxito de las herramientas); executor se ejecuta bajo un guardrail de seguridad. Esta arquitectura de tres capas le permite al agent mantener coherencia en tareas de contexto largo que abarcan varias herramientas.
\end{importantbox}

\subsection{Tool Orchestration y Workflow}

Para que el Agent coordine varias herramientas (búsqueda, bases de datos, API), se necesita un flujo de trabajo y una tool specification claros:

\begin{itemize}
    \item Definir tool intent, inputs y outputs, para evitar model hallucinating tool usage
    \item Usar `tool descriptor` para validar los parámetros de llamada del LLM (tipo, rango)
    \item Insertar `pre-conditions` y `post-conditions` en el workflow
\end{itemize}

\begin{knowledgebox}{Diseño del workflow como state machine}
Un agent complejo a menudo modela el workflow como una state machine: cada estado representa una fase de la tarea (comprensión, planning, execution, confirmation); la condición que dispara la transition proviene del juicio del planner. Esto permite simular distintas ramas con unit test durante la fase de desarrollo, y mejora la system reliability.
\end{knowledgebox}
\begin{knowledgebox}{El valor del tool descriptor}
Cuando el modelo necesita invocar `fetch\_legislation()` o `run\_cad\_simulation()`, el tool descriptor ofrece schema, authority y expected latency, para evitar que el agent llame a ciegas o que se agote el tiempo de espera.
\end{knowledgebox}

\subsection{Resumen de la sección}

Una plataforma de IA empresarial necesita ofrecer flexibilidad de elección en cada una de las capas: chips, modelos, herramientas y aplicaciones.

